\section{Interfaces, reflection and energy partition}\label{sec:interfaces}
Reflection and transmission at interfaces is the most heavily studied wave problem in generalized thermoelasticity and multiphysics media; Table~\ref{tab:streams} lists many of the works. Most of them report \emph{energy ratios} and use the requirement that the ratios sum to unity as a validation. This section builds the problem from the elastic baseline, then adds thermal coupling and shows when and why the unit-sum check is satisfied or violated, since this is a point on which the literature is not uniform.

\subsection{Welded elastic interface at normal incidence}
Two elastic half-spaces with impedances $Z_j=\rho_jc_{1j}$ are perfectly bonded at $x=0$. For a wave incident from medium~1, continuity of displacement and of traction gives the displacement amplitude coefficients
\begin{equation}
r=\frac{Z_1-Z_2}{Z_1+Z_2},\qquad t=\frac{2Z_1}{Z_1+Z_2},\label{eq:rt}
\end{equation}
and the time-averaged energy fractions $\mathcal R=r^2$, $\mathcal T=4Z_1Z_2/(Z_1+Z_2)^2$ satisfy $\mathcal R+\mathcal T=1$ (\Cref{fig:interface}). The displacement transmission coefficient $t$ can reach two for a soft second medium without violating conservation, since only the energy fractions are bounded.
\begin{figure}[t]\centering
\includegraphics[width=\textwidth]{figs/fig_interface.pdf}
\caption{Normal incidence on a welded elastic interface. (a) Displacement reflection and transmission coefficients. (b) Energy fractions; $\mathcal R+\mathcal T=1$ for every impedance ratio.}\label{fig:interface}
\end{figure}

\subsection{Oblique incidence on an elastic interface}
For oblique incidence of a P wave, Snell's law conserves the horizontal slowness $p=\sin\theta_1/c_{1,P}=\sin\phi_1/c_{1,S}=\sin\theta_2/c_{2,P}=\sin\phi_2/c_{2,S}$ and four waves are generated (reflected P and SV, transmitted P and SV). Continuity of $u_x,u_z,\sigma_{zz},\sigma_{xz}$ gives four linear equations for the amplitudes (the Zoeppritz system). We solved it for a pair of media ($\rho_1,c_P,c_S=2.5,5.0,2.8$ and $\rho_2,c_P,c_S=3.0,8.0,4.5$ in arbitrary consistent units), with energy flux per wave $P_z=-\tfrac12\Real(\sigma_{zj}\overline{\dot u_j})$. The sum of the four energy fractions equals one to rounding error ($|\text{sum}-1|<2\times10^{-15}$ over $0.5^\circ$--$89.5^\circ$), and at normal incidence the P--P reflection fraction ($0.099268$) equals the impedance formula ($0.099268$). \Cref{fig:ienergy}(a) shows the characteristic features: a critical angle ($38.7^\circ$ here), beyond which the transmitted P wave is evanescent and carries no energy; a rise of the reflected P fraction to total reflection of P energy in the P--P channel only at grazing incidence, with mode conversion in between.

\subsection{Thermoelastic interfaces}
When both media are thermoelastic, two further wave types appear in each medium and the interface conditions include thermal continuity. For perfect thermal contact, $\theta$ and the normal flux $q_n$ are continuous together with displacement and traction. In normal incidence there are four amplitudes (reflected elastic and thermal waves, transmitted elastic and thermal waves) and four equations. We implemented this for two media with different densities, wave speeds, heat capacities, conductivities, coupling constants and lag times, with Lord--Shulman conduction; incident wave of elastic type. The amplitude ratios and the flux of each wave follow from \eqref{eq:P}. The tests we performed are: (i) as $\varepsilon\to0$ (here $\beta=10^{-6}$) the reflected elastic energy fraction ($0.169550173$) equals the impedance formula $(Z_1-Z_2)^2/(Z_1+Z_2)^2$ to twelve digits; (ii) the \emph{net} flux, computed from the total fields on either side of the interface, is continuous to $10^{-15}$ for all frequencies tested.

\paragraph{When does ``energy fractions sum to one'' hold?} The exact statement is item (ii): the net flux of \eqref{eq:P} into the interface equals the net flux out of it, because the continuity conditions imply continuity of $\sigma\dot u$ and of $\theta q$. This is not the same as saying that the \emph{sum of the fluxes carried by each wave separately} equals the incident flux. For homogeneous waves in lossless media the two statements coincide, because the cross terms between distinct modes vanish in the time average (this is what the elastic checks above confirm numerically). In a dissipative medium, the modes have complex wavenumbers and the cross terms between, for example, the elastic-type and the thermal-type wave at the interface do not vanish, so
\begin{equation}
\sum_n\mathcal F_n\ne1\quad\text{in general, while}\quad \mathcal F_{\rm net}\equiv1.
\end{equation}
\Cref{fig:ienergy}(b,c) shows an example: the sum of the four mode fractions is within $1.5\,\%$ of unity for $\omega\lesssim2\omega_*$ but reaches 1.34 near $\omega\approx8\omega_*$ and individual fractions can exceed one (the transmitted elastic fraction reaches 1.1). This is not a violation of energy conservation: the cross terms account for the difference. The consequence for published results is practical. An energy-ratio sum that deviates from unity can be caused by (a) a wrong coupling coefficient or boundary condition, (b) a flux definition inconsistent with the model, or (c) genuine interference of modes in a dissipative medium; the first two are errors, the third is not, and a correct check must use the net flux, or include the interference terms, rather than the sum of mode fluxes. Studies that report energy ratios at thermoelastic and piezothermoelastic interfaces \citep{Barak2023,Barak2023IJP,Barak2023JTS,Barak2023MMMS,Gupta2023Energy,Kumar2023PhysScr,Kumar2019,Kumar2022,Kumar2024,Kumar2024a} should be read with this point in mind: when the media are lossy, a unit sum of mode ratios is not a necessary condition and a failure of the sum to equal one is not by itself evidence of error.

\begin{figure}[t]\centering
\includegraphics[width=\textwidth]{figs/fig_interface_energy.pdf}
\caption{Energy partition at welded interfaces. (a) Oblique P incidence on an elastic--elastic interface: reflected and transmitted P and SV energy fractions; their sum is unity at every angle. (b) Normal incidence on a thermoelastic--thermoelastic interface with Lord--Shulman conduction ($\varepsilon_1=0.1$, $\varepsilon_2=0.3$, $\tau_{q1}=0.25/\omega_*$, $\tau_{q2}=0.15/\omega_*$): energy fractions of the four generated waves against frequency. (c) The sum of mode fractions deviates from unity at high frequency because of interference terms, while the net flux of \eqref{eq:P} is continuous across the interface (dashed).}\label{fig:ienergy}
\end{figure}

\subsection{Imperfect interfaces}
Real interfaces are rarely perfect. The common idealisation replaces continuity by spring conditions: the traction is continuous and proportional to the jump in displacement, $\sigma=K_n[u]$ (normal) and $K_t[u_t]$ (tangential); thermal imperfections add a thermal contact resistance, $[\theta]=R_\theta q$ (the Kapitza-type jump), and for micropolar or piezoelectric media analogous electric and rotational conditions. In the limit $K\to\infty$ one recovers the welded interface and for $K\to0$ a free surface. With thermal resistance the net flux of \eqref{eq:P} is no longer continuous: the jump $[\theta]q/T_0=R_\theta q^2/T_0$ is the entropy production at the interface, and the energy balance must include it. Imperfect-interface wave problems have been treated for micropolar solids \citep{Fu2013}, piezoelectric composites \citep{Kumawat2023,Gupta2025}, and a fluid--double-porous thermoelastic interface \citep{Kumar2025}.

\subsection{Fluid--solid and multiphysics interfaces}
When one medium is an ideal fluid the shear conditions change (no shear traction, tangential displacement discontinuous), and the energy ratios have been computed for fluid--piezothermoelastic and fluid--double-porous interfaces \citep{Gupta2023IJNMF,Kumar2019,Kumar2025}. In piezoelectric media electric boundary conditions (continuity of potential and of normal electric displacement, or an electrode condition) add two equations per interface; the reflection and refraction at elastic--magneto-thermoelastic diffusion solids with voids \citep{KumarKansal2017}, micropolar piezothermoelastic half-spaces \citep{Kumar2018}, functionally graded piezoelectric microstructures \citep{Liu2020} and flexoelectric or flexomagnetic microstructures \citep{Singh2020c,Jiao2024} extend the same procedure. All of these share the structure of Sections~\ref{sec:surface} and the present one: a linear system for amplitudes whose entries are fixed by the dispersion roots of each medium, and an energy flux that must be defined consistently with the model.

\subsection{A checklist for interface calculations}
The following checks, which cost little, catch most errors in interface calculations:
\begin{enumerate}
\item Reduce the coupling to zero and compare with the elastic impedance formulas (normal incidence) or the Zoeppritz equations.
\item Verify the Snell invariants (horizontal slowness) for all generated waves, including evanescent ones.
\item Verify continuity of the net flux \eqref{eq:P} (not of the sum of mode fluxes) across a perfect interface; for imperfect interfaces include the interface dissipation.
\item Test the boundary conditions independently: substitute the computed amplitudes back into each boundary equation.
\item For lagged and memory models, state the flux definition used for the energy ratio and whether interference terms are included.
\end{enumerate}
